\قسمت{ترتیب حرکت با کلید \متن‌لاتین{Tab}}

در فرم ثبت‌نامی که نوشته‌اید، کاربر پس از پر کردن «نام» کلید \متن‌لاتین{Tab} را می‌زند و به جای رفتن به «نام خانوادگی»، مکان‌نما روی دکمه‌ی «ثبت» در انتهای فرم می‌رود.
دلیلش این است که برای چیدمان، ترتیب المان‌ها در \متن‌لاتین{HTML} با آنچه کاربر می‌بیند یکی نیست.

\شروع{شمارش}
\فقره با کدام صفت\پانویس{Attribute} می‌توان این ترتیب را اصلاح کرد؟
\فقره چرا دادن مقدارهای بزرگ‌تر از صفر به آن (مثلا ۱، ۲، ۳ روی هر ورودی) راه‌حل خوبی نیست؟
\پایان{شمارش}

\نمره{۱}

% پاسخ پیشنهادی است و پیش از استفاده در تصحیح نیاز به بازبینی دارد.
\begin{پاسخ}

صفت \متن‌لاتین{tabindex}.

مقدار بزرگ‌تر از صفر یک ترتیب جداگانه می‌سازد که \متن‌سیاه{پیش از} همه‌ی المان‌های عادی صفحه پیمایش می‌شود؛
با اضافه شدن هر المان تازه باید همه‌ی شماره‌ها را بازبینی کرد و به‌سادگی از هماهنگی خارج می‌شود.
راه درست، مرتب کردن خود \متن‌لاتین{HTML} به ترتیب منطقی و استفاده از \متن‌لاتین{tabindex="0"} در صورت نیاز است.

\شروع{فقرات}
\فقره نام بردن \متن‌لاتین{tabindex}: ۵۰ درصد نمره
\فقره بیان مشکل مقدارهای مثبت یا پیشنهاد اصلاح ترتیب \متن‌لاتین{HTML}: ۵۰ درصد نمره
\پایان{فقرات}

\end{پاسخ}
